\section{Playoff Success Factors}
\label{sec:playoff-factors}

\subsection{The ``Defense Wins Championships'' Literature}
\label{sec:dwc-literature}

The phrase ``defense wins championships'' occupies a peculiar position in sports discourse: it is perhaps the most widely cited empirical claim in professional football, yet it has been subjected to remarkably little rigorous academic scrutiny. The saying functions simultaneously as folk wisdom (invoked by commentators and coaches without citation), as strategic prescription (informing roster construction and draft philosophy), and as empirical hypothesis (making a testable claim about the relationship between defensive quality and postseason success). Its persuasive power derives partly from this ambiguity: by remaining vague about the precise mechanism, magnitude, and conditions under which defense ``wins'' championships, the claim resists straightforward falsification while maintaining rhetorical force.

\textcite{robst2011defense} published what remains the most direct peer-reviewed examination of the hypothesis in American football. Using data from 1981 through 2009, they employed logistic regression to predict Super Bowl participation and victory from a battery of traditional performance metrics: total offensive yards, total defensive yards allowed, points scored, points allowed, turnover differential, and strength of schedule. Their central finding was that offensive metrics  -- particularly points scored and offensive yards  -- were substantially stronger predictors of Super Bowl appearance and victory than defensive metrics. They concluded that the ``defense wins championships'' maxim lacks empirical support in the modern NFL, at least when evaluated using traditional box-score statistics.

The limitations of their approach, however, constrain the inferential weight that can be placed on this conclusion. Their independent variables  -- total yards and total points  -- are aggregate measures that conflate quality with opportunity and context. A team that plays with a lead for much of every game will allow additional yards and points in garbage time that inflate its defensive statistics without reflecting genuine defensive weakness. Conversely, an excellent defense on a bad team may face additional possessions because its own offense goes three-and-out frequently, accumulating opportunities for the opponent to gain yards regardless of per-drive defensive quality. These confounds are well understood in the modern analytics community but were not addressed in Robst et al.'s methodology. Furthermore, their sample ended in 2009  -- before the full impact of the NFL's rule changes favoring passing offenses, and before the data infrastructure existed to support play-level or drive-level analysis.

No subsequent peer-reviewed study has directly revisited the hypothesis with modern analytical tools. The claim has been debated extensively in the grey literature  -- sports analytics blogs, podcasts, and social media  -- but these discussions typically rely on selective examples (citing recent champions with elite defenses while ignoring those with elite offenses), small-sample anecdotes (a single Super Bowl matchup treated as evidence for a general rule), or metrics with the reproducibility problems documented in \Cref{sec:dvoa,sec:pff}. The present thesis fills this gap by bringing calibrated machine learning, explicit decomposition, and multi-method hypothesis testing to a question that has been debated through anecdote for decades.

\subsection{Quantitative Studies of Postseason Determinants}
\label{sec:postseason-determinants}

Beyond the specific ``defense wins championships'' debate, a broader literature has examined what factors predict playoff success in professional sports \parencite{baker2013assessing}. This research spans multiple leagues and methodological traditions, providing context for the present investigation.

In baseball, the sabermetric revolution demonstrated convincingly that run prevention (pitching and defense) and run production (hitting) contribute approximately equally to regular-season winning, a finding formalized in the Pythagorean winning percentage formula \parencite{miller2007pythagorean}. However, the playoff prediction literature in baseball reveals that regular-season quality only weakly predicts postseason success due to the high variance inherent in short playoff series \parencite{lopez2018randomness}. This ``crapshoot'' characterization of baseball playoffs has been documented extensively and suggests caution in any sport when attributing postseason outcomes to specific quality dimensions: when sample sizes are small (a single-elimination game or a best-of-three series), randomness dominates systematic quality differences \parencite{boulier2003predicting}.

The NFL postseason amplifies this concern because every playoff game is single-elimination. A team plays at most four postseason games en route to a Super Bowl victory, meaning that a single turnover, injury, or officiating decision can end a season regardless of the team's overall quality \parencite{lopez2018randomness}. This high variance creates a fundamental tension for any analysis of ``what wins championships'': if randomness is a dominant factor in single-game outcomes, then the relationship between any quality dimension and championship success will necessarily be weak, and large samples are required to detect genuine effects. The present paper addresses this through multi-season aggregation (224 team-seasons across 2018--2024) and multiple statistical methods that can detect effects of varying magnitudes.

\textcite{massey2013loserscurse} approached the question from a labor market perspective, examining whether NFL teams make rational decisions about talent acquisition. Their finding that teams systematically overvalue early draft picks  -- paying a premium for high-round selections that does not correspond to proportionally higher player production  -- implies a market-level inefficiency in talent evaluation. While their study focused on individual player value rather than team-level offense-defense allocation, their methodology and findings are relevant to the present investigation in two ways. First, they demonstrated that NFL decision-making is subject to systematic cognitive biases rather than being perfectly rational, which is consistent with the hypothesis that the ``defense wins championships'' heuristic persists not because it is empirically correct but because it is culturally entrenched. Second, their use of surplus value analysis  -- comparing player production to draft-pick cost  -- provides a conceptual template for the resource allocation analysis discussed in \Cref{sec:resource-allocation}.

\subsection{Home-Field Advantage and Contextual Factors}
\label{sec:home-field}

The playoff success literature has consistently identified home-field advantage as a significant predictor of postseason outcomes \parencite{courneya1992home}. Teams playing at home win at elevated rates in the NFL playoffs, with estimated effects ranging from 55--60\% historical win rates for home teams compared to the 50\% null expectation \parencite{boulier2003predicting}. This advantage arises from a combination of factors: crowd noise disrupting opposing offensive communication, elimination of travel fatigue, familiarity with playing surface and weather conditions, and potential officiating bias \parencite{courneya1992home}.

Home-field advantage interacts with the offense-defense question in a subtle way. Higher-seeded teams earn home-field advantage by performing well in the regular season, and the quality dimension that drives regular-season success therefore determines which teams enjoy the postseason home-field benefit. If offensive quality is a stronger determinant of regular-season wins (a finding consistently supported by the EPA literature and by the results of this paper), then offensively dominant teams disproportionately earn home-field advantage, which compounds their playoff success probability. The observed association between offensive dominance and playoff advancement may therefore reflect both a direct effect (better offenses win more games) and an indirect effect mediated through seeding (better offenses earn home-field advantage, which provides an additional win-probability boost). The present analysis does not attempt to decompose these direct and indirect effects, which would require a structural equation modeling approach beyond the scope of this paper, but acknowledges their theoretical relevance.

\subsection{Era Effects and Rule Changes}
\label{sec:era-effects}

Any investigation of championship predictors in the NFL must contend with the reality that the competitive environment has shifted considerably over time. The league has implemented numerous rule changes since 2004 that have collectively shifted the balance of the game toward passing offenses and away from physical, coverage-based defensive play. The most consequential include:

\begin{itemize}
  \item The 2004 ``Ty Law Rule'' (officially: illegal contact enforcement), which restricted defensive backs from initiating contact with receivers beyond five yards from the line of scrimmage, opening downfield passing lanes that defenses had previously closed through physicality \parencite{nflops2004rules}.

  \item Expanded roughing-the-passer protections (implemented iteratively from 2006 through 2018), which penalized defenders for hits to the quarterback's head and neck area, landing on the quarterback with body weight, and hitting below the knee \parencite{pereira2018roughing}. These protections extended the effective time quarterbacks have to throw and reduced the physical deterrent that pass rushes previously imposed.

  \item The 2010 emphasis on helmet-to-helmet hit enforcement, which penalized defenders for targeting receivers in vulnerable positions, reducing the physical cost of catching passes over the middle of the field.

  \item The 2019 expansion of pass interference reviewability (adopted following a controversial non-call in the 2018 NFC Championship Game) and its non-renewal for 2020, which generated league-wide awareness of defensive contact and may have influenced defensive coaching toward less aggressive coverage techniques.
\end{itemize}

The cumulative effect of these changes is visible in league-wide passing efficiency trends. Passer ratings, completion percentages, yards per attempt, and passing touchdowns have all risen markedly since 2004, while defensive metrics such as interception rates and sack rates have declined \parencite{barnwell2019passingrevolution}. The 2020s represent the most offense-friendly era in NFL history by virtually every statistical measure \parencite{nflscoring2023trends}.

These era effects have two implications for the present investigation. First, they create a temporal confound: the ``defense wins championships'' maxim may have been correct in the run-heavy, low-scoring NFL of the 1970s and 1980s but incorrect in the modern passing-dominated game. The present paper deliberately restricts its sample to 2018--2024, a period within the modern passing era, so that its findings reflect contemporary rather than historical competitive dynamics. This choice sacrifices sample size for temporal relevance: the results speak specifically to whether defensive dominance predicts championship success in the current NFL, not whether it ever did so in the past.

Second, the rule changes suggest a mechanistic explanation for why offensive quality might dominate in the modern game. If the rules systematically constrain defensive tactics (limiting contact, protecting quarterbacks, penalizing aggressive coverage) while leaving offensive innovation unrestricted (spread formations, RPO concepts, pre-snap motion), then the ceiling on defensive impact is artificially lowered while the ceiling on offensive impact remains high. A team can build a historically great offense through scheme innovation and talent acquisition; the same investment in defense faces diminishing returns because the rules prevent the defense from fully executing the tactics that elite defensive talent enables. This mechanism would explain not only why offensive quality predicts championships in the modern era, but why the relationship may have been weaker or reversed in earlier periods when defensive physicality was less constrained.
